% Table 1: Main results across all datasets
\begin{table}[t]
\centering
\caption{Same-model performance on Qwen3-4B across 5 datasets. \textbf{PAIR+SESS} and \textbf{DAN+SESS} denote PAIR iterative optimization augmented with SESS skill retrieval. DAN+SESS uses 6 fixed DAN templates as cold-start initialization (structural prior), not AutoDAN's genetic evolution.}
\label{tab:main_results}
\resizebox{\columnwidth}{!}{%
\small
\begin{tabular}{lccccc}
\toprule
\textbf{Dataset} & \textbf{PAIR} & \textbf{PAIR+SESS} & \textbf{$\Delta$} & \textbf{AutoDAN} & \textbf{DAN+SESS} \\
\midrule
default & 55.5 & \textbf{79.0} & \textbf{+23.5} & 86.8 & \textbf{100.0} \\
advbench & 77.3 & \textbf{81.0} & +3.7 & --- & \textbf{100.0} \\
harmbench\_ctx & 81.0 & \textbf{98.0} & \textbf{+17.0} & --- & \textbf{100.0} \\
harmbench\_std & 75.5 & \textbf{91.5} & \textbf{+16.0} & --- & \textbf{100.0} \\
jailbreakBench & 72.0 & \textbf{88.0} & \textbf{+16.0} & --- & \textbf{100.0} \\
\midrule
\textbf{Average} & \textbf{72.3} & \textbf{87.5} & \textbf{+15.2} & \textbf{86.8} & \textbf{100.0} \\
\bottomrule
\end{tabular}%
}
\vspace{0.3em}
\footnotesize
\textit{Note:} ``---" indicates AutoDAN baseline not evaluated on that dataset. DAN+SESS retains PAIR's attacker-target dialogue loop with structural priors; AutoDAN evolves prompts via genetic operations without iterative refinement.
\end{table}

% 说明：
% 1. PAIR baseline 55.5% 来自 transfer/Qwen3-4B/default/pair
% 2. AutoDAN baseline 86.8% 来自 transfer/Qwen3-4B/default/autodan
% 3. PAIR+SESS 79.0% 来自 transfer/Qwen3-4B/default/pair_skills_28
% 4. DAN+SESS 100.0% 来自 transfer/Qwen3-4B/default/autodan_skills_1
% 5. DAN+SESS 使用 6 个固定 DAN 模板作为初始 skills（结构先验），而非 AutoDAN 的进化机制
